\documentclass[11pt,a4paper]{article}

\usepackage[T1]{fontenc}
\usepackage[utf8]{inputenc}
%\usepackage{lmodern}
\usepackage[margin=25mm]{geometry}
\usepackage{amsmath,amssymb}
\usepackage{siunitx}
\usepackage{booktabs}
\usepackage{graphicx}
\usepackage[siunitx,RPvoltages]{circuitikz}
\usepackage{tikz}
\usepackage{hyperref}
\usepackage{caption}

\hypersetup{colorlinks=true,linkcolor=black,citecolor=black,urlcolor=blue}
\sisetup{per-mode=symbol}

\title{The Scott-T Transformer:\\Three-Phase to Two-Phase Conversion\\and its Simscape Implementation}
\author{}
\date{}

\newcommand{\ph}[1]{\underline{#1}}

\begin{document}
\maketitle

\begin{abstract}
This note describes the Scott connection, which converts a three-phase system
into a two-phase system whose two voltages have equal magnitude and are
displaced by \SI{90}{\degree}. The connection rule and the
$\sqrt{3}/2$ turns ratio of the teaser winding are derived, the current
relations and the symmetrical-component behaviour under unbalanced two-phase
load are given, and the equations implemented in the accompanying Simscape
component \texttt{scott\_t.ssc} are documented together with its per-unit
parameterisation and sign conventions.
\end{abstract}

\section{Introduction}

A three-phase to two-phase transformation is required whenever a load or a
source that is intrinsically two-phase, or a pair of single-phase loads that
must be seen as balanced by the supply, has to be interfaced to a three-phase
network. Typical applications are AC railway substations feeding two adjacent
single-phase catenary sections, induction furnaces and electric arc furnaces,
two-phase servo and control motors, and laboratory supplies for two-phase test
benches.

The classical solution is the Scott connection, introduced by C.~F.~Scott at
the end of the nineteenth century and still described in standard transformer
texts \cite{fitzgerald,winders}. It uses two single-phase transformers, called
the \emph{main} and the \emph{teaser}, and produces two secondary voltages of
equal magnitude in exact quadrature. A functionally equivalent alternative, the
Le~Blanc connection, achieves the same result on a single three-limb core at
the cost of a more complex winding arrangement \cite{chenguo}.

The answer to the question that motivated this note is therefore affirmative:
the two secondary phases of a Scott-connected transformer are displaced by
exactly \SI{90}{\degree}, and this is a property of the connection geometry, not
an approximation.

\section{Connection}

Figure~\ref{fig:conn} shows the arrangement. The main transformer primary,
with $N_1$ turns, is connected between lines A and B and is centre-tapped at
node~O. The teaser transformer primary, with $N_{1t}$ turns, is connected
between line C and the centre tap~O. Both secondaries have $N_2$ turns.

\begin{figure}[htbp]
\centering
\begin{circuitikz}[scale=1.0,transform shape,american]

% --- main transformer -------------------------------------------------
\draw (0,4.2) node[circle,fill,inner sep=1.2pt]{} node[above]{A}
      to[L=$N_1/2$] (0,2.6) -- (0,2.1) coordinate (O);
\draw (O) -- (0,1.6) to[L=$N_1/2$] (0,0) node[circle,fill,inner sep=1.2pt]{}
      node[below]{B};
\draw (O) node[circle,fill,inner sep=1.2pt]{} node[right=1pt]{O};
% main core
\draw[very thick] (0.95,-0.2) -- (0.95,4.4);
\draw[very thick] (1.15,-0.2) -- (1.15,4.4);
% main secondary
\draw (1.7,4.2) node[circle,fill,inner sep=1.2pt]{} node[above]{$a_1$}
      to[L=$N_2$] (1.7,0) node[circle,fill,inner sep=1.2pt]{}
      node[below]{$a_2$};
\node at (2.9,0.6) {\small main};

% --- teaser transformer -----------------------------------------------
\draw (-6.2,2.1) node[circle,fill,inner sep=1.2pt]{} node[left]{C}
      to[L={$N_{1t}=\tfrac{\sqrt{3}}{2}N_1$}] (-3.4,2.1) -- (O);
% teaser core
\draw[very thick] (-6.4,1.55) -- (-3.2,1.55);
\draw[very thick] (-6.4,1.35) -- (-3.2,1.35);
% teaser secondary
\draw (-6.2,0.4) node[circle,fill,inner sep=1.2pt]{} node[left]{$b_1$}
      to[L=$N_2$] (-3.4,0.4) node[circle,fill,inner sep=1.2pt]{}
      node[right]{$b_2$};
\node at (-4.8,3.0) {\small teaser};

\end{circuitikz}
\caption{Scott connection. The teaser primary is returned to the centre tap of
the main primary, so that its current splits equally between the two halves of
the main winding and produces no net magneto-motive force in the main core.}
\label{fig:conn}
\end{figure}

The essential feature is the centre tap. The teaser current returns through the
two halves of the main primary in opposite winding senses, so the two
contributions to the main-core magneto-motive force cancel. The two magnetic
circuits are therefore decoupled even though the electrical circuits are not.

\section{Voltage relations and the teaser turns ratio}

Let the primary be supplied by a balanced positive-sequence system of phase
voltages
\begin{equation}
\ph{V}_A = V\!\angle 0^\circ,\quad
\ph{V}_B = V\!\angle\!-\!120^\circ,\quad
\ph{V}_C = V\!\angle 120^\circ .
\end{equation}
Neglecting winding impedances, the voltage impressed on the main primary is
\begin{equation}
\ph{V}_{AB} = \ph{V}_A - \ph{V}_B = \sqrt{3}\,V\!\angle 30^\circ .
\end{equation}
Because the two halves of the main primary are identical and carry the same
induced voltage, the centre tap sits at the mean of the two line potentials,
\begin{equation}
\ph{V}_O = \tfrac{1}{2}\left(\ph{V}_A + \ph{V}_B\right) = -\tfrac{1}{2}\ph{V}_C ,
\end{equation}
where the last equality uses $\ph{V}_A+\ph{V}_B+\ph{V}_C = 0$. The voltage
impressed on the teaser primary is therefore
\begin{equation}
\ph{V}_{CO} = \ph{V}_C - \ph{V}_O = \tfrac{3}{2}\,\ph{V}_C
            = \tfrac{3}{2}\,V\!\angle 120^\circ .
\end{equation}

Two facts follow immediately. First, $\ph{V}_{CO}$ leads $\ph{V}_{AB}$ by
exactly \SI{90}{\degree}: the quadrature is geometric and holds for any balanced
supply. Second, the two primary voltages are \emph{not} equal in magnitude,
their ratio being $\tfrac{3}{2}V / \sqrt{3}V = \sqrt{3}/2$. To obtain equal
secondary voltages with equal secondary turns, the teaser primary must have
\begin{equation}
\boxed{\;N_{1t} = \frac{\sqrt{3}}{2}\,N_1 \approx 0.8660\,N_1\;}
\end{equation}
Denoting $n = N_1/N_2$ and $V_1 = \sqrt{3}V$ the primary line-to-line rms
voltage, the secondary voltages are
\begin{align}
\ph{V}_a &= \frac{\ph{V}_{AB}}{n} = \frac{V_1}{n}\angle 30^\circ, &
\ph{V}_b &= \frac{\ph{V}_{CO}}{N_{1t}/N_2} = \frac{V_1}{n}\angle 120^\circ .
\end{align}
The rated secondary phase voltage is thus $V_2 = V_1/n$ and $\ph{V}_b$ leads
$\ph{V}_a$ by \SI{90}{\degree}. The phasor construction is shown in
Fig.~\ref{fig:phasor}.

\begin{figure}[htbp]
\centering
\begin{tikzpicture}[scale=1.5,>=stealth]
% primary phasors
\draw[->,thick] (0,0) -- (30:1.732) node[right]{$\ph{V}_{AB}$};
\draw[->,thick] (0,0) -- (120:1.5)  node[above left]{$\ph{V}_{CO}$};
\draw[->,gray]  (0,0) -- (0:1)      node[below right]{$\ph{V}_{A}$};
\draw[->,gray]  (0,0) -- (-120:1)   node[below]{$\ph{V}_{B}$};
\draw[->,gray]  (0,0) -- (120:1)    node[left]{$\ph{V}_{C}$};
% right angle mark
\draw (30:0.42) arc (30:120:0.42);
\node at (75:0.60) {\footnotesize $90^\circ$};
\draw[dashed] (0,0) -- (30:0.42);
\end{tikzpicture}
\caption{Phasor diagram. The main primary sees the line-to-line voltage
$\ph{V}_{AB}$, the teaser primary sees $\tfrac{3}{2}\ph{V}_{C}$; the two are in
quadrature and their magnitudes are in the ratio $2/\sqrt{3}$.}
\label{fig:phasor}
\end{figure}

\section{Current relations}

Let $\ph{I}_a$ and $\ph{I}_b$ be the currents delivered by the main and teaser
secondaries and let magnetizing currents be neglected. Magneto-motive force
balance on the two cores gives
\begin{align}
\frac{N_1}{2}\left(\ph{I}_{1a} + \ph{I}_{1b}\right) &= N_2 \ph{I}_a, &
N_{1t}\,\ph{I}_C &= N_2 \ph{I}_b ,
\end{align}
where $\ph{I}_{1a}$ is the current from A to O and $\ph{I}_{1b}$ the current
from O to B, so that $\ph{I}_{1a} = \ph{I}_A$ and $\ph{I}_{1b} = -\ph{I}_B$.
With Kirchhoff's current law at node~O, $\ph{I}_A+\ph{I}_B+\ph{I}_C = 0$, the
line currents are
\begin{align}
\ph{I}_A &= \frac{1}{n}\left(\ph{I}_a - \tfrac{1}{\sqrt{3}}\ph{I}_b\right), &
\ph{I}_B &= \frac{1}{n}\left(-\ph{I}_a - \tfrac{1}{\sqrt{3}}\ph{I}_b\right), &
\ph{I}_C &= \frac{2}{\sqrt{3}\,n}\,\ph{I}_b .
\label{eq:lines}
\end{align}
For a balanced two-phase load, $\ph{I}_b = \mathrm{j}\ph{I}_a$, and
\eqref{eq:lines} reduces to a balanced three-phase set of magnitude
$2\lvert\ph{I}_a\rvert/(\sqrt{3}n)$.

\subsection*{Behaviour under unbalanced two-phase load}

Applying the symmetrical-component transformation to \eqref{eq:lines} yields
the compact result
\begin{align}
\ph{I}_{+} &= \frac{e^{-\mathrm{j}30^\circ}}{\sqrt{3}\,n}
              \left(\ph{I}_a - \mathrm{j}\ph{I}_b\right), &
\ph{I}_{-} &= \frac{e^{+\mathrm{j}30^\circ}}{\sqrt{3}\,n}
              \left(\ph{I}_a + \mathrm{j}\ph{I}_b\right).
\label{eq:seq}
\end{align}
The negative-sequence current vanishes if and only if
$\ph{I}_b = \mathrm{j}\ph{I}_a$, that is, only for a two-phase load that is
balanced both in magnitude and in phase angle. The Scott connection does not
by itself compensate an unbalanced load; it only maps a balanced two-phase load
onto a balanced three-phase supply. In traction substations this is the reason
why the two catenary sections must be loaded symmetrically, or why additional
compensation is required \cite{chenguo}.

\subsection*{Rating}

Each transformer carries one half of the total apparent power, $S_r/2$. The
main transformer is rated for $V_1$ on the primary and the teaser for
$(\sqrt{3}/2)V_1$, both for the same secondary voltage $V_2$. The two units are
therefore not interchangeable. In practice the two windings are often made
identical up to the tap position, and both are provided with the
$\sqrt{3}/2 \approx 0.866$ tap plus the centre tap, so that either unit can act
as main or as teaser.

\section{Simscape model}

The component \texttt{scott\_t.ssc} implements the connection as a single
Simscape block with seven external electrical nodes: \texttt{A}, \texttt{B},
\texttt{C} on the primary side and \texttt{a1}, \texttt{a2}, \texttt{b1},
\texttt{b2} on the secondary side. The centre tap \texttt{O} is an internal
private node. The two secondaries are galvanically separate; tying
\texttt{a2} to \texttt{b2} externally produces a two-phase three-wire system
with a common return.

\subsection{Equations}

Winding resistances and leakage inductances are lumped in series with each
winding, and each core is represented by a linear magnetizing inductance in
parallel with a core-loss resistance, both referred to the secondary. Denoting
by $e_m$ and $e_t$ the secondary induced electromotive forces, the implemented
equations are

\begin{align}
v_A - v_O &= R_{1h}\,i_{1a} + L_{1h}\,\frac{\mathrm{d}i_{1a}}{\mathrm{d}t}
             + \frac{n}{2}\,e_m ,\\[2pt]
v_O - v_B &= R_{1h}\,i_{1b} + L_{1h}\,\frac{\mathrm{d}i_{1b}}{\mathrm{d}t}
             + \frac{n}{2}\,e_m ,\\[2pt]
v_C - v_O &= R_{1t}\,i_{1t} + L_{1t}\,\frac{\mathrm{d}i_{1t}}{\mathrm{d}t}
             + n_t\,e_t ,\\[2pt]
v_{a1} - v_{a2} &= e_m - R_2\,i_{2m}
             - L_2\,\frac{\mathrm{d}i_{2m}}{\mathrm{d}t} ,\\[2pt]
v_{b1} - v_{b2} &= e_t - R_2\,i_{2t}
             - L_2\,\frac{\mathrm{d}i_{2t}}{\mathrm{d}t} ,
\end{align}
with the magneto-motive force balance referred to the secondary windings,
\begin{align}
\frac{n}{2}\left(i_{1a} + i_{1b}\right) - i_{2m}
   &= i_{Lm} + \frac{e_m}{R_{fe}} , &
n_t\,i_{1t} - i_{2t} &= i_{Lt} + \frac{e_t}{R_{fe}} ,
\end{align}
and the magnetizing branches
\begin{align}
e_m &= L_m\,\frac{\mathrm{d}i_{Lm}}{\mathrm{d}t} , &
e_t &= L_m\,\frac{\mathrm{d}i_{Lt}}{\mathrm{d}t} .
\end{align}
Here $n = N_1/N_2 = V_1/V_2$ and $n_t = N_{1t}/N_2 = (\sqrt{3}/2)\,n$.
Kirchhoff's current law at the internal node,
$i_{1a} + i_{1t} = i_{1b}$, is enforced automatically by the branch
declarations.

Two points are worth noting. First, the main primary is split into two halves
of $N_1/2$ turns each, with half the total resistance and half the total
leakage inductance; the two halves are not mutually coupled through their
leakage paths, which is the usual approximation. Second, the sum
$i_{1a}+i_{1b}$ drives the main core while their difference, equal to the
teaser return current, does not. The centre-tap cancellation is thus a property
of the equations and does not need to be imposed separately.

\subsection{Parameters}

The impedances are entered in per unit. Each transformer processes $S_r/2$, so
the secondary base impedance is
\begin{equation}
Z_{b2} = \frac{V_2^2}{S_r/2} = \frac{2V_2^2}{S_r} ,
\end{equation}
and the primary bases follow from the respective turns ratios,
$Z_{b1,\text{main}} = n^2 Z_{b2}$ and
$Z_{b1,\text{teaser}} = n_t^2 Z_{b2} = \tfrac{3}{4}n^2 Z_{b2}$. With this
choice the same per-unit value applies to both units. Table~\ref{tab:par}
lists the parameters.

\begin{table}[htbp]
\centering
\caption{Parameters of the \texttt{scott\_t} component.}
\label{tab:par}
\begin{tabular}{llll}
\toprule
Name & Description & Unit & Default \\
\midrule
\texttt{Sr}  & Rated apparent power, three-phase total & \si{\volt\ampere} & \num{100e3} \\
\texttt{V1}  & Primary rated line-to-line voltage, rms & \si{\volt} & 400 \\
\texttt{V2}  & Secondary rated phase voltage, rms      & \si{\volt} & 230 \\
\texttt{fr}  & Rated frequency                         & \si{\hertz} & 50 \\
\texttt{r1}  & Primary winding resistance              & pu & 0.005 \\
\texttt{l1}  & Primary leakage inductance              & pu & 0.030 \\
\texttt{r2}  & Secondary winding resistance            & pu & 0.005 \\
\texttt{l2}  & Secondary leakage inductance            & pu & 0.030 \\
\texttt{lm}  & Magnetizing inductance                  & pu & 100 \\
\texttt{rfe} & Core-loss resistance                    & pu & 500 \\
\bottomrule
\end{tabular}
\end{table}

\subsection{Sign conventions}

Positive current enters the component at \texttt{A}, \texttt{B} and \texttt{C}
and leaves it at \texttt{a1} and \texttt{b1}. With a balanced positive-sequence
supply of line-to-line rms voltage \texttt{V1} applied with phase sequence
A--B--C, the open-circuit secondary voltages are
\begin{align}
v_{a1}-v_{a2} &= \sqrt{2}\,V_2\cos(\omega t + 30^\circ), &
v_{b1}-v_{b2} &= \sqrt{2}\,V_2\cos(\omega t + 120^\circ),
\end{align}
so the teaser secondary leads the main secondary by \SI{90}{\degree}. Reversing
the phase sequence at the primary reverses the sign of the quadrature, which is
the expected behaviour and a useful check.

\subsection{Limitations}

The magnetic model is linear: saturation, hysteresis and frequency-dependent
core losses are not represented, and no thermal port is provided. Inrush
phenomena are therefore not reproduced. Leakage coupling between the two halves
of the main primary is neglected, which slightly underestimates the impedance
seen by the teaser return current. If saturation is required, the two
magnetizing relations $e = L_m\,\mathrm{d}i_L/\mathrm{d}t$ can be replaced by a
flux-linkage formulation $e = \mathrm{d}\lambda/\mathrm{d}t$ with a nonlinear
$\lambda(i_L)$ characteristic, at the cost of one additional differential
variable per core.

\section{Suggested verification}

A minimal test model consists of a balanced three-phase voltage source
connected to \texttt{A}, \texttt{B}, \texttt{C}, with \texttt{a2} tied to
\texttt{b2} and a resistive load on each secondary. Three checks are
sufficient:
\begin{enumerate}
\item \emph{No load.} The two secondary voltages must be equal in rms value and
in exact quadrature, and the magnetizing currents must be in the ratio
$1:1$ referred to the secondary.
\item \emph{Balanced load.} With equal resistive loads on both secondaries, the
three line currents must form a balanced positive-sequence set; by
\eqref{eq:seq} the negative-sequence component must be zero to within the
magnetizing current.
\item \emph{Single-phase load.} Loading only the main secondary must produce a
negative-sequence line current of magnitude equal to the positive-sequence one,
in agreement with \eqref{eq:seq} for $\ph{I}_b = 0$.
\end{enumerate}

\begin{thebibliography}{9}

\bibitem{fitzgerald}
A.~E.~Fitzgerald, C.~Kingsley~Jr., and S.~D.~Umans,
\emph{Electric Machinery}, 6th ed. New York, NY, USA: McGraw-Hill, 2003.

\bibitem{winders}
J.~J.~Winders~Jr.,
\emph{Power Transformer Principles and Applications}.
New York, NY, USA: Marcel Dekker, 2002.

\bibitem{chenguo}
B.-K.~Chen and B.-S.~Guo,
``Three-phase models of specially connected transformers,''
\emph{IEEE Trans. Power Del.}, vol.~11, no.~1, pp.~323--330, Jan. 1996.

\bibitem{ieee}
\emph{IEEE Standard for General Requirements for Liquid-Immersed Distribution,
Power, and Regulating Transformers}, IEEE Std C57.12.00.

\bibitem{simscape}
The MathWorks, Inc., \emph{Simscape Language Guide}, Natick, MA, USA.

\end{thebibliography}

\end{document}
